% Wersja źródłowa: audyt II, jednostka 15A
\setcounter{section}{15}
\noindent\textbf{Wymagane wyniki.} Korzystamy z twierdzenia Hahna--Banacha, twierdzenia Banacha--Steinhausa, podstawowego twierdzenia rachunku dla krzywych o wartościach w przestrzeni Banacha i faktu, że ciągłe funkcjonały liniowe rozdzielają punkty lokalnie wypukłej przestrzeni Hausdorffa. Pełną teorię całek na takich przestrzeniach rozwijamy poniżej.
\medskip
% \hyperref[sec:integral-lcs15a]{Moduł 15A}. Celowo nie zmieniamy numeracji \hyperref[sec:wariacyjny16]{rozdziałów 16}--33.
\begingroup
\renewcommand{\thesection}{15A}
\renewcommand{\theHsection}{15A}
\section{Całkowanie funkcji o wartościach w przestrzeniach lokalnie wypukłych}
\label{sec:integral-lcs15a}

\providecommand{\atlasref}[2]{#2}


W \hyperref[sec:gateaux-frechet14]{rozdziale~14} udowodniliśmy podstawowe twierdzenie
Newtona--Leibniza dla krzywych w przestrzeni Banacha, potrzebne w rachunku
różniczkowym. Teraz budujemy szerszą teorię: obiektem całkowanym może być
funkcja przyjmująca wartości w przestrzeni lokalnie wypukłej, której topologia
nie musi pochodzić od jednej normy. Nie należy utożsamiać tego zagadnienia
z całkowaniem funkcji mierzalnych: całki Bochnera, Pettisa, Dunforda i Gelfanda
będą przedmiotem \hyperref[sec:integral-vector15b]{modułu~15B}.

W całym module $\K\in\{\R,\C\}$, a $E$ jest przestrzenią liniowo-topologiczną
Hausdorffa nad $\K$, lokalnie wypukłą, z topologią wyznaczoną przez rodzinę
ciągłych seminorm $\mathcal P$. \emph{Zupełność ciągowa} oznacza, że każdy
ciąg Cauchy'ego w $E$ ma granicę w $E$; nie mylimy jej z ogólną zupełnością
względem sieci Cauchy'ego. Jeśli nie zaznaczono inaczej, $a<b$ są liczbami
rzeczywistymi. Funkcja na odcinku jest rozumiana z topologią podprzestrzeni.

\subsection{Sumy Riemanna i twierdzenie o istnieniu}

\noindent\emph{Powiązanie tematyczne w Bibliotece matematycznej:} \atlasref{AM1-29::ma1:def:riemann-sum}{AM1-29: sumy Riemanna dla funkcji skalarnych}.


\begin{definicja}[Całka Riemanna o wartościach w przestrzeni liniowo-topologicznej]
\label{def:riemann-lcs15a}
Niech $a<b$ będą liczbami rzeczywistymi, $\K\in\{\R,\C\}$,
$E$ przestrzenią lokalnie wypukłą Hausdorffa nad $\K$,
$\mathcal P$ rodziną ciągłych seminorm wyznaczających jej
topologię i $f:[a,b]\to E$ funkcją. \emph{Podziałem z punktami pośrednimi} nazywamy
parę $P=(a=t_0<\cdots<t_N=b;\,\xi_j\in[t_{j-1},t_j])$.
Piszemy $|P|=\max_j(t_j-t_{j-1})$ i
$S(f,P)=\sum_{j=1}^N(t_j-t_{j-1})f(\xi_j)\in E$,
gdzie $N\ge1$ i $1\le j\le N$ w warunku na punkty pośrednie.
Funkcja $f$ jest całkowalna w sensie Riemanna w $E$, jeżeli istnieje $I\in E$
takie, że dla każdej seminormy $p\in\mathcal P$ i $\varepsilon>0$ istnieje
$\delta>0$ taka, iż $p(S(f,P)-I)<\varepsilon$ dla \emph{każdego} podziału
z punktami pośrednimi spełniającego $|P|<\delta$.
Warunek jest równoważny zbieżności sum Riemanna
w topologii $E$, ponieważ skończone przecięcia kul seminorm
tworzą bazę otoczeń zera. W szczególności zbieżność zachodzi
względem każdej ciągłej seminormy, nie tylko seminorm z $\mathcal P$.
Granica jest jednoznaczna w przestrzeni Hausdorffa.
Wówczas oznaczamy $I=\int_a^b f(t)\,\mathrm dt$.
Dla funkcji określonej na odcinku o końcach $a,b$ przy
$a=b$ przyjmujemy całkę równą zeru, a przy $b<a$ definiujemy
$\int_a^b f(t)\,\mathrm dt=-\int_b^a f(t)\,\mathrm dt$,
jeżeli całka po prawej stronie istnieje.
\end{definicja}

\begin{uwaga}[Znaczenie topologii]
Definicja wymaga zbieżności \emph{w topologii $E$}, a nie tylko po
zastosowaniu każdego funkcjonału $\ell\in E'$. Tę drugą własność nazywamy
słabą całkowalnością skalarną i sama nie gwarantuje istnienia wektora całki.
Powyższa granica jest granicą rodziny wszystkich sum Riemanna przy
$|P|\to0$; nie sprowadzamy jej z definicji do jednego wybranego ciągu podziałów.
\end{uwaga}

\begin{twierdzenie}[Istnienie całki przy zupełności ciągowej]
\label{thm:riemann-existence15a}
Niech $a<b$ będą liczbami rzeczywistymi, $\K\in\{\R,\C\}$,
a $E$ zupełną ciągowo przestrzenią lokalnie wypukłą Hausdorffa
nad $\K$. Wtedy każda funkcja ciągła $f:[a,b]\to E$
jest całkowalna w sensie Riemanna. W szczególności wynik
zachodzi w każdej zupełnej przestrzeni lokalnie wypukłej Hausdorffa.
\end{twierdzenie}
\begin{proof}
Dla $p\in\mathcal P$ definiujemy moduł ciągłości
$\omega_p(r)=\sup\{p(f(t)-f(s)):s,t\in[a,b],\ |s-t|\le r\}$.
Ciągłość $f$ na zwartym odcinku oznacza jednostajną ciągłość
w topologii $E$, zatem $\omega_p(r)\to0$ dla $r\downarrow0$.

Niech $P,Q$ będą dwoma podziałami z punktami pośrednimi, o średnicach
co najwyżej $r$. Rozważmy wspólne rozdrobnienie ich podziałów
przedziału i wybierzmy na nim dowolne punkty pośrednie, otrzymując $R$.
Każdy przedział $R$ mieści się w jednym przedziale $P$, a odpowiadające
mu punkty pośrednie są odległe o najwyżej $r$; analogicznie dla $Q$.
Po rozpisaniu sum względem wspólnego rozdrobnienia dostajemy
\[
 p(S(f,P)-S(f,R))\le(b-a)\omega_p(r),\qquad
 p(S(f,Q)-S(f,R))\le(b-a)\omega_p(r).
\]
Zatem $p(S(f,P)-S(f,Q))\le2(b-a)\omega_p(r)$.

Wybierzmy ciąg podziałów dyadycznych $P_n$, na przykład ze środkami
przedziałów jako punktami pośrednimi, o średnicach dążących do zera.
Powyższa nierówność pokazuje, że $(S(f,P_n))$ jest ciągiem Cauchy'ego
w $E$. Z zupełności ciągowej istnieje $I\in E$, do którego zbiega.
Dla dowolnego $P$ o średnicy mniejszej od $r$ stosujemy ostatnie
oszacowanie do $P$ i $P_n$ przy dostatecznie dużym $n$, a następnie
przechodzimy do granicy $n\to\infty$:
$p(S(f,P)-I)\le2(b-a)\omega_p(r)$.
Ponieważ dotyczy to każdej seminormy, zachodzi warunek z definicji.
Jeżeli $I,J\in E$ spełniają warunek z definicji, to dla
każdego $p\in\mathcal P$ i $\varepsilon>0$ wybieramy podział
$P$ tak drobny, że $p(I-S(f,P))<\varepsilon/2$ i
$p(S(f,P)-J)<\varepsilon/2$. Z nierówności trójkąta
$p(I-J)<\varepsilon$, więc $p(I-J)=0$. Ponieważ $\mathcal P$
rozdziela punkty, mamy $I=J$.
\end{proof}

\begin{stwierdzenie}[Liniowość i oszacowania seminormowe]
\label{prop:riemann-seminorm15a}
Niech $a<b$ będą rzeczywiste, $\K\in\{\R,\C\}$,
$E$ przestrzenią lokalnie wypukłą Hausdorffa nad $\K$,
a $f,g:[a,b]\to E$ funkcjami ciągłymi i całkowalnymi w $E$.
Dla $\alpha,\beta\in\K$ funkcja $\alpha f+\beta g$ jest
całkowalna i zachodzi pierwsza z poniższych równości.
Druga zachodzi dla $c\in[a,b]$, jeśli ograniczenia $f$ do $[a,c]$
i $[c,b]$ są całkowalne w $E$ (w szczególności gdy $E$ jest zupełna ciągowo):
\[
 \int_a^b(\alpha f(t)+\beta g(t))\,\mathrm dt
 =\alpha\int_a^b f(t)\,\mathrm dt+\beta\int_a^b g(t)\,\mathrm dt,
 \qquad \int_a^b f(t)\,\mathrm dt=\int_a^c f(t)\,\mathrm dt
 +\int_c^b f(t)\,\mathrm dt.
\]
Dla stałej funkcji o wartości $v\in E$ zachodzi
$\int_a^b v\,\mathrm dt=(b-a)v$.
Jeżeli $E$ jest zupełna ciągowo, to dla wszystkich
$u,v,w\in[a,b]$ zachodzi ponadto addytywność całek z orientacją:
$\int_u^v f(t)\,\mathrm dt+\int_v^w f(t)\,\mathrm dt
=\int_u^w f(t)\,\mathrm dt$.
Dla każdej ciągłej seminormy $p$ na $E$ mamy
\[
 p\left(\int_a^b f(t)\,\mathrm dt\right)
 \le\int_a^b p(f(t))\,\mathrm dt
 \le (b-a)\sup_{t\in[a,b]}p(f(t)).
\]
\end{stwierdzenie}
\begin{proof}
Oznaczmy całki $f,g$ przez $I_f,I_g$. Z liniowości sum
Riemanna otrzymujemy
$p(S(\alpha f+\beta g,P)-(\alpha I_f+\beta I_g))
\le|\alpha|p(S(f,P)-I_f)+|\beta|p(S(g,P)-I_g)$.
Dla każdej seminormy $p$ prawa strona dąży do zera
wraz ze średnicą podziału, co daje istnienie całki
kombinacji liniowej i wskazaną równość.
Jeżeli $a<c<b$, wybieramy ciągi podziałów odcinków $[a,c]$
i $[c,b]$ o średnicach dążących do zera. Ich połączenia są
podziałami $[a,b]$ o średnicach dążących do zera, a odpowiadające
sumy Riemanna są sumami dwóch sum na pododcinkach.
Przechodząc do granicy, uzyskujemy addytywność.
Dla $c=a$ lub $c=b$ równość wynika z konwencji dla całki
na odcinku zdegenerowanym. Dla funkcji stałej każda suma
Riemanna wynosi $(b-a)v$, co dowodzi wzoru.
Przy zupełności ciągowej z twierdzenia o istnieniu mamy
całki na wszystkich pododcinkach. Oznaczmy
$H(x)=\int_a^x f(t)\,\mathrm dt$. Addytywność dla
$a\le u\le v\le b$ daje $\int_u^v f(t)\,\mathrm dt=H(v)-H(u)$.
Na mocy konwencji orientacji tożsamość zachodzi również
dla $v<u$ oraz $u=v$. Dla dowolnych $u,v,w$ sumujemy
dwie takie różnice, otrzymując żądaną addytywność.
Nierówność trójkąta i dodatnia jednorodność seminormy dają
$p(S(f,P))\le\sum_j(t_j-t_{j-1})p(f(\xi_j))$.
Po przejściu do granicy lewa strona zbiega do
$p(\int_a^b f(t)\,\mathrm dt)$, a prawa do skalarnej
całki funkcji ciągłej $p\circ f$. Ostatnia nierówność jest skalarna.
\end{proof}

\begin{twierdzenie}[Przemienność z odwzorowaniami liniowymi ciągłymi]
\label{thm:riemann-functorial15a}
Niech $a<b$ będą rzeczywiste, $\K\in\{\R,\C\}$,
$E,F$ przestrzeniami lokalnie wypukłymi Hausdorffa nad $\K$,
$f:[a,b]\to E$ funkcją ciągłą i całkowalną w $E$, a
$T:E\to F$ odwzorowaniem liniowym ciągłym. Wtedy $T\circ f$
jest całkowalna w $F$ i
\[
 T\left(\int_a^b f(t)\,\mathrm dt\right)
 =\int_a^b T(f(t))\,\mathrm dt.
\]
W szczególności równość zachodzi dla każdego $\ell\in E'$, a
całkę w $E$ wyznaczają jednoznacznie jej wartości skalarne
$\ell(\int_a^b f(t)\,\mathrm dt)$ dla $\ell\in E'$.
Nie zakładamy zupełności $F$.
\end{twierdzenie}
\begin{proof}
Dla każdego $P$ mamy $S(T\circ f,P)=T(S(f,P))$.
Ciągłość $T$ przenosi granicę sum na $T(I)$; ten wektor jest
całką w $F$, nawet jeśli $F$ nie jest zupełna.
Funkcjonały z $E'$ rozdzielają punkty $E$ na mocy
twierdzenia Hahna--Banacha i hausdorffowskości $E$,
więc równości wszystkich obrazów skalarnych wyznaczają $I$.
\end{proof}

\begin{wniosek}[Środek ciężkości krzywej]
\label{cor:barycenter15a}
Niech $a<b$ będą rzeczywiste, $\K\in\{\R,\C\}$,
$E$ zupełną ciągowo przestrzenią lokalnie wypukłą Hausdorffa
nad $\K$, a $f:[a,b]\to E$ funkcją ciągłą. Wówczas
\[
 (b-a)^{-1}\int_a^b f(t)\,\mathrm dt
 \in\overline{\operatorname{conv}}\bigl(f([a,b])\bigr),
\]
gdzie domknięcie i otoczkę wypukłą rozumiemy w topologii $E$.
\end{wniosek}
\begin{proof}
Dla każdego podziału $P$ wektor $(b-a)^{-1}S(f,P)$ jest
kombinacją wypukłą wartości $f(\xi_j)$, ponieważ jego
współczynniki $(t_j-t_{j-1})/(b-a)$ są nieujemne i sumują się do $1$.
Zamkniętość wskazanego zbioru i zbieżność sum dowodzą tezy.
\end{proof}

\subsection{Rola zupełności i całka niewłaściwa}

\begin{przyklad}[Funkcja ciągła bez całki w przestrzeni niezupełnej]
\label{ex:noncomplete-integral15a}
Rozważmy $E=c_{00}(\N;\R)$, przestrzeń ciągów rzeczywistych
o skończonym nośniku z normą odziedziczoną z $\ell^2(\N;\R)$.
Niech $e_n$ oznacza $n$-ty wektor jednostkowy.
Obcięcia każdego ciągu z $\ell^2$ do pierwszych $N$ współrzędnych
zbiegają w normie $\ell^2$ do tego ciągu, ponieważ ogon
szeregu kwadratów modułów dąży do zera. Stąd $c_{00}$ jest
gęsta w $\ell^2$, lecz nie jest domknięta (np. obcięcia
$(2^{-n})_{n\ge1}$ zbiegają poza $c_{00}$).
Jest zatem niezupełna, a jej uzupełnieniem jest $\ell^2$.
Wybierzmy nieujemną, niezerową funkcję ciągłą $\eta:\R\to\R$
o zwartym nośniku zawartym w $(0,1)$.
Dla $t\in[0,1]$ niech
\[
 f(t)=\sum_{n=1}^{\infty}2^{-n}\eta(2^{n+1}t-1)e_n.
\]
Przedziały, na których kolejne składniki są niezerowe,
leżą w rozłącznych zbiorach $(2^{-(n+1)},2^{-n})$.
Dlatego dla każdego $t$ suma ma co najwyżej jeden niezerowy
wyraz i $f(t)\in c_{00}$. W otoczeniu każdego $t_0>0$ tylko skończenie wiele składników
może być niezerowych, gdyż prawe końce ich przedziałów dążą
do zera. Funkcja jest tam skończoną sumą funkcji ciągłych.
Mamy $f(0)=0$; dla $0<t<2^{-N}$ jedyny ewentualny
niezerowy składnik ma indeks $n\ge N$, więc
$\|f(t)\|_2\le2^{-N}\|\eta\|_\infty\to0$ przy $N\to\infty$.
To dowodzi ciągłości w zerze.
Całka istnieje w $\ell^2$, lecz jej $n$-ta współrzędna to
\[
 \int_0^1 2^{-n}\eta(2^{n+1}t-1)\,\mathrm dt
 =2^{-2n-1}\int_0^1\eta(s)\,\mathrm ds>0.
\]
Ma ona więc nieskończenie wiele niezerowych współrzędnych i nie
należy do $c_{00}$. Gdyby $f$ była całkowalna w $E$, przemienność
z ciągłymi funkcjonałami współrzędnych wymusiłaby tę samą
wartość o nieskończonym nośniku; sprzeczność.

Co więcej, definiując $g:[0,2]\to c_{00}$ przez $g(t)=f(t)$ dla
$t\in[0,1]$ oraz $g(t)=-f(2-t)$ dla $t\in[1,2]$, otrzymujemy
funkcję ciągłą, ponieważ $f(1)=0$. Jej całka w $\ell^2$ jest zerem: dla każdego ciągłego
funkcjonału $\ell$ podstawienie $s=2-t$ daje
$\int_1^2\ell(g(t))\,\mathrm dt=-\int_0^1\ell(f(s))\,\mathrm ds$.
Przemienność całek z funkcjonałami i rozdzielanie punktów
dają tę samą równość całek wektorowych, a addytywność
daje całkę zero. Sumy Riemanna $g$ zbiegają
więc do zera także w normie odziedziczonej przez $c_{00}$.
Jednak ograniczenie $g$ do $[0,1]$ nie jest całkowalne w $c_{00}$.
Pokazuje to konieczność dodatkowego założenia w addytywności całki
względem przedziału.
\end{przyklad}

\begin{stwierdzenie}[Całka w uzupełnieniu]
\label{prop:completion-integral15a}
Niech $a<b$ będą rzeczywiste, $\K\in\{\R,\C\}$,
$E$ przestrzenią lokalnie wypukłą Hausdorffa nad $\K$,
a $j:E\to\widehat E$ jej kanonicznym zanurzeniem do uzupełnienia.
Dla ciągłej $f:[a,b]\to E$ funkcja $j\circ f$ jest
całkowalna w $\widehat E$. Całka $f$ w $E$ istnieje wtedy
i tylko wtedy, gdy $\int_a^b j(f(t))\,\mathrm dt\in j(E)$;
wówczas $j(\int_a^b f(t)\,\mathrm dt)
=\int_a^b j(f(t))\,\mathrm dt$.
\end{stwierdzenie}
\begin{proof}
Uzupełnienie $\widehat E$ jest zupełne, zatem pierwsza teza
wynika z twierdzenia o istnieniu. Jeżeli całka w $E$ istnieje,
przemienność z $j$ daje konieczność. W drugą stronę załóżmy,
że całka w $\widehat E$ jest postaci $j(v)$ dla $v\in E$.
Rodzina sum Riemanna należy do $j(E)$ i zbiega w $\widehat E$
do $j(v)$. Topologia $E$ jest topologią indukowaną przez
zanurzenie $j$, więc sumy te zbiegają do $v$ w $E$.
\end{proof}

\begin{twierdzenie}[Kryterium całki niewłaściwej]
\label{thm:improper-lcs15a}
Niech $a\in\R$, $\K\in\{\R,\C\}$, $E$ będzie przestrzenią
lokalnie wypukłą Hausdorffa nad $\K$, zupełną ciągowo,
a $f:[a,\infty)\to E$ funkcją ciągłą.
Załóżmy, że dla każdej ciągłej seminormy $p$ zachodzi
$\int_a^\infty p(f(t))\,\mathrm dt<\infty$ w sensie
niewłaściwej całki skalarnej. Wówczas istnieje
$\lim_{R\to\infty}\int_a^R f(t)\,\mathrm dt\in E$, oznaczana przez
$\int_a^\infty f(t)\,\mathrm dt$, oraz
$p(\int_a^\infty f(t)\,\mathrm dt)
\le\int_a^\infty p(f(t))\,\mathrm dt$.
\end{twierdzenie}
\begin{proof}
Dla $R'>R\ge a$ z oszacowania seminormowego i addytywności
całki wynika
$p(\int_a^{R'}f(t)\,\mathrm dt-\int_a^R f(t)\,\mathrm dt)
\le\int_R^{R'}p(f(t))\,\mathrm dt$.
Ogon ostatniej całki dąży do zera, dlatego ciąg
$\bigl(\int_a^{a+n}f(t)\,\mathrm dt\bigr)_{n\ge1}$
jest ciągiem Cauchy'ego w $E$ i ma granicę $I$.
Dla ustalonych $p$ i $\varepsilon>0$ wybieramy $n$ tak duże,
że $p(\int_a^{a+n}f(t)\,\mathrm dt-I)<\varepsilon/2$
oraz $\int_{a+n}^\infty p(f(t))\,\mathrm dt<\varepsilon/2$.
Dla każdego rzeczywistego $R\ge a+n$ mamy zatem
$p(\int_a^R f(t)\,\mathrm dt-I)<\varepsilon$ z nierówności
trójkąta i oszacowania ogona. To daje granicę dla rzeczywistego
parametru $R\to\infty$. Dla każdego $R$ zachodzi
$p(\int_a^R f(t)\,\mathrm dt)\le\int_a^R p(f(t))\,\mathrm dt$;
ciągłość seminormy i granica całki skalarnej dają oszacowanie końcowe.
\end{proof}

\subsection{Twierdzenie Newtona--Leibniza i rachunek z parametrem}

\begin{twierdzenie}[Podstawowe twierdzenie rachunku całkowego]
\label{thm:ftc-lcs15a}
Niech $a<b$ będą rzeczywiste, a $E$ zupełną ciągowo
przestrzenią lokalnie wypukłą Hausdorffa nad $\K\in\{\R,\C\}$.
Krzywą klasy $C^1$ na $[a,b]$ rozumiemy jako krzywą ciągłą
z pochodną w topologii $E$ wewnątrz przedziału i pochodnymi
jednostronnymi na końcach, tworzącymi ciągłą funkcję
na $[a,b]$. Niech $f:[a,b]\to E$ będzie ciągła. Wtedy funkcja
$F:[a,b]\to E$, $F(x)=\int_a^x f(t)\,\mathrm dt$, jest klasy
$C^1$ na $[a,b]$ w powyższym sensie oraz
$F'(x)=f(x)$ dla każdego $x\in[a,b]$ (z odpowiednimi
pochodnymi jednostronnymi).
Jeżeli $G:[a,b]\to E$ jest klasy $C^1$, to
\[
 G(b)-G(a)=\int_a^b G'(t)\,\mathrm dt.
\]
W szczególności jest to uogólnienie wyniku z \hyperref[sec:gateaux-frechet14]{rozdziału~14},
a nie jego ponowny dowód w przypadku przestrzeni Banacha.
\end{twierdzenie}
\begin{proof}
Dla $x,x+h\in[a,b]$ i $h\ne0$ z addytywności całki mamy
$h^{-1}(F(x+h)-F(x))-f(x)
=h^{-1}\int_x^{x+h}(f(t)-f(x))\,\mathrm dt$.
Oszacowanie seminormowe (również dla $h<0$ po zmianie
orientacji całki) daje
\[
 p\left(h^{-1}(F(x+h)-F(x))-f(x)\right)
 \le\sup_{\substack{t\in[a,b]\\|t-x|\le|h|}}p(f(t)-f(x))
 \longrightarrow0.
\]
Stąd pochodna istnieje w topologii $E$ i jest równa $f(x)$.
Istnienie skończonej pochodnej (również jednostronnej)
zapewnia ciągłość $F$ w $x$: przyrost jest iloczynem $h$
i ilorazu różnicowego zbieżnego do $f(x)$.
Pochodna jest ciągła na $[a,b]$, bo $f$ jest ciągła.

Dla drugiej tezy nie musimy odtwarzać dowodu z przestrzeni Banacha.
Dla każdego $\ell\in E'$ złożenie $\ell\circ G$ jest skalarne
klasy $C^1$ z pochodną $\ell(G'(t))$. Skalarne twierdzenie
Newtona--Leibniza i przemienność całki z $\ell$ dają
$\ell(G(b)-G(a))=\ell(\int_a^b G'(t)\,\mathrm dt)$.
Ponieważ $E'$ rozdziela punkty, zachodzi żądana równość w $E$.
\end{proof}

\begin{twierdzenie}[Różniczkowanie całki zależnej od parametru]
\label{thm:parametric-integral15a}
Niech $a<b$ będą rzeczywiste, $\K\in\{\R,\C\}$,
$E$ zupełną ciągowo przestrzenią lokalnie wypukłą Hausdorffa
nad $\K$, $J\subseteq\R$ otwartym przedziałem i
$f:[a,b]\times J\to E$ funkcją ciągłą. Załóżmy, że dla
każdego $(t,s)\in[a,b]\times J$ istnieje pochodna
$\partial_s f(t,s)$ w topologii $E$ i że funkcja
$(t,s)\mapsto\partial_s f(t,s)$ jest ciągła.
Wtedy $F(s)=\int_a^b f(t,s)\,\mathrm dt$ jest funkcją klasy $C^1$
o wartościach w $E$ oraz
$F'(s)=\int_a^b\partial_s f(t,s)\,\mathrm dt$ dla $s\in J$.
\end{twierdzenie}
\begin{proof}
Ustalmy $s\in J$ i wybierzmy $\delta>0$ tak, aby
$[s-\delta,s+\delta]\subset J$. Dla każdej ciągłej seminormy $p$
mamy $p(F(s+h)-F(s))\le(b-a)\sup_{t\in[a,b]}
p(f(t,s+h)-f(t,s))\to0$ dzięki jednostajnej ciągłości
na zwartym prostokącie. Dla $0<|h|<\delta$ stosujemy
wektorowe twierdzenie Newtona--Leibniza do krzywej
$u\mapsto f(t,u)$, otrzymując
$f(t,s+h)-f(t,s)=\int_s^{s+h}\partial_s f(t,u)\,\mathrm du$.
Po odjęciu $h\partial_s f(t,s)$, podzieleniu przez $h$
i użyciu oszacowania seminormowego dostajemy
\[
 p\left(\frac{f(t,s+h)-f(t,s)}h-\partial_s f(t,s)\right)
 \le\sup_{u\text{ między }s\text{ i }s+h}
 p(\partial_s f(t,u)-\partial_s f(t,s)).
\]
Z liniowości całki i oszacowania seminormowego mamy
\[
 p\left(\frac{F(s+h)-F(s)}h
       -\int_a^b\partial_s f(t,s)\,\mathrm dt\right)
 \le(b-a)\sup_{\substack{t\in[a,b]\\u\text{ między }s\text{ i }s+h}}
 p(\partial_s f(t,u)-\partial_s f(t,s)).
\]
Jednostajna ciągłość $\partial_s f$ na wybranym zwartym
prostokącie daje zbieżność prawej strony do zera, więc
pochodna ma wskazaną postać. Dla jej ciągłości stosujemy
$p(F'(s+h)-F'(s))\le(b-a)\sup_{t\in[a,b]}
p(\partial_s f(t,s+h)-\partial_s f(t,s))\to0$.
\end{proof}


\begin{twierdzenie}[Reguła Leibniza z ruchomymi granicami]
\label{thm:leibniz-moving15a}
Niech $\K\in\{\R,\C\}$, $E$ będzie zupełną ciągowo
przestrzenią lokalnie wypukłą Hausdorffa nad $\K$,
$I,J$ otwartymi przedziałami rzeczywistymi, $f:I\times J\to E$ ciągłą funkcją
z pochodną $\partial_s f$ istniejącą w topologii $E$
w każdym punkcie i ciągłą na $I\times J$, a $a,b:J\to I$
funkcjami klasy $C^1$. Całkę z odwróconymi granicami
rozumiemy jako ujemną całkę w porządku rosnącym.
Wtedy $F(s)=\int_{a(s)}^{b(s)}f(t,s)\,\mathrm dt$
jest klasy $C^1$ w topologii $E$ i
\[
 F'(s)=f(b(s),s)b'(s)-f(a(s),s)a'(s)
       +\int_{a(s)}^{b(s)}\partial_s f(t,s)\,\mathrm dt.
\]
\end{twierdzenie}
\begin{proof}
Ustalmy $s\in J$. Wybieramy zwarty przedział $K\subset I$
zawierający $a(s),b(s)$ we wnętrzu i $\delta>0$ tak, że
$[s-\delta,s+\delta]\subset J$ i dla $|h|\le\delta$ obie
granice $a(s+h),b(s+h)$ należą do $K$.
Jest to możliwe dzięki ciągłości $a,b$ i otwartości $I,J$. Wszystkie całki
istnieją; addytywność całek z orientacją daje
\[
 \begin{split}
 F(s+h)-F(s)={}&
 \int_{a(s)}^{b(s)}(f(t,s+h)-f(t,s))\,\mathrm dt\\
 &+\int_{b(s)}^{b(s+h)} f(t,s+h)\,\mathrm dt
 -\int_{a(s)}^{a(s+h)} f(t,s+h)\,\mathrm dt.
 \end{split}
\]
Pierwszy składnik po podzieleniu przez $h$ dąży
do całki pochodnej na mocy twierdzenia o stałych
granicach (po odwróceniu granic zmieniamy znak,
a przy $a(s)=b(s)$ obie całki są zerowe).
Dla dowolnej ciągłej seminormy $p$
różnica drugiego ilorazu od
$((b(s+h)-b(s))/h)f(b(s),s)$ ma seminormę co najwyżej
\[
 \left|\frac{b(s+h)-b(s)}h\right|
 \sup_{t\text{ między }b(s),b(s+h)}
 p(f(t,s+h)-f(b(s),s)).
\]
Pierwszy czynnik pozostaje ograniczony, a drugi
dąży do zera z ciągłości $f$ w punkcie $(b(s),s)$.
Analogiczne oszacowanie dla $a$ daje wzór pochodnej,
również gdy granice są równe lub zmieniają kolejność.
Ciągłość odwzorowań $u\mapsto f(b(u),u)b'(u)$
i $u\mapsto f(a(u),u)a'(u)$ wynika z ciągłości $f,a,b,a',b'$
i ciągłości mnożenia przez skalar w $E$. Dla ostatniego oznaczmy $g(t,u)=\partial_s f(t,u)$ oraz
$G(u)=\int_{a(u)}^{b(u)}g(t,u)\,\mathrm dt$.
Rozbicie przyrostu jak wyżej i oszacowanie seminormowe dają
\[
\begin{split}
 p(G(s+h)-G(s))\le{}&
 |b(s)-a(s)|\sup_{t\in K}p(g(t,s+h)-g(t,s))\\
 &+\bigl(|b(s+h)-b(s)|+|a(s+h)-a(s)|\bigr)M_p,
\end{split}
\]
gdzie $M_p=\sup_{(t,u)\in K\times[s-\delta,s+\delta]}p(g(t,u))$
jest skończone przez zwartość. Pierwszy składnik dąży
do zera dzięki jednostajnej ciągłości w seminormie $p$,
drugi dzięki ciągłości $a,b$. To daje ciągłość $G$ w $E$.
\end{proof}

\begin{twierdzenie}[Wielokrotne różniczkowanie całki z parametrem]
\label{thm:higher-parametric15a}
Niech $a<b$ będą rzeczywiste, $m\ge1$, $\K\in\{\R,\C\}$,
$E$ zupełną ciągowo przestrzenią lokalnie wypukłą Hausdorffa
nad $\K$ i $J\subseteq\R$ otwartym przedziałem. Jeżeli $f:[a,b]\times J\to E$ jest ciągła, a każda
pochodna $\partial_s^j f$, $1\le j\le m$, istnieje
w topologii $E$ i jest wspólnie ciągła względem $(t,s)$,
to
$F(s)=\int_a^b f(t,s)\,\mathrm dt$ jest klasy $C^m$ i
$F^{(j)}(s)=\int_a^b\partial_s^j f(t,s)\,\mathrm dt$
dla $1\le j\le m$.
\end{twierdzenie}
\begin{proof}
Przypadek $j=1$ jest twierdzeniem o różniczkowaniu
całki z parametrem. Jeżeli wzór zachodzi dla $j<m$,
stosujemy to samo twierdzenie do funkcji $\partial_s^j f$,
której pochodna względem $s$ jest ciągła na mocy
założeń. Otrzymujemy wzór dla $j+1$ i ciągłość tej
pochodnej. Indukcja kończy dowód.
\end{proof}


\begin{twierdzenie}[Twierdzenie Fubiniego dla funkcji ciągłych]
\label{thm:fubini-continuous15a}
Niech $a<b$ i $c<d$ będą liczbami rzeczywistymi,
$\K\in\{\R,\C\}$, $E$ zupełną ciągowo przestrzenią
lokalnie wypukłą Hausdorffa nad $\K$ i
$f:[a,b]\times[c,d]\to E$ funkcją ciągłą. Wówczas obie funkcje
$s\mapsto\int_c^d f(s,t)\,\mathrm dt$ oraz
$t\mapsto\int_a^b f(s,t)\,\mathrm ds$ są ciągłe i
\[
 \int_a^b\left(\int_c^d f(s,t)\,\mathrm dt\right)\mathrm ds
 =\int_c^d\left(\int_a^b f(s,t)\,\mathrm ds\right)\mathrm dt.
\]
Nie zakładamy mierzalności Bochnerowskiej ani twierdzenia Fubiniego
dla ogólnych funkcji mierzalnych.
\end{twierdzenie}
\begin{proof}
Dla każdej seminormy $p$ jednostajna ciągłość $f$
na zwartym prostokącie daje ciągłość obydwu funkcji
całek iterowanych: na przykład
$p(\int_c^d(f(s,t)-f(s_0,t))\,\mathrm dt)
\le(d-c)\sup_{t\in[c,d]}p(f(s,t)-f(s_0,t))$.

Niech $P=(s_j;\xi_j)$ i $Q=(t_k;\eta_k)$ będą podziałami
z punktami pośrednimi odpowiednio odcinków $[a,b]$ i $[c,d]$;
piszemy $\Delta s_j=s_j-s_{j-1}$ i $\Delta t_k=t_k-t_{k-1}$.
Dla $p$ oznaczmy przez $\omega_p^{(1)}(r)$ i $\omega_p^{(2)}(r)$
suprema seminormy różnicy wartości $f$ przy zmianie tylko
odpowiedniej zmiennej o co najwyżej $r$.
Oba dążą do zera dla $r\downarrow0$ dzięki jednostajnej
ciągłości na prostokącie. Niech
$D(P,Q)=\sum_{j,k}\Delta s_j\Delta t_k f(\xi_j,\eta_k)$.
Dla $A(s)=\int_c^d f(s,t)\,\mathrm dt$ mamy
\[
\begin{split}
 p(S(A,P)-D(P,Q))
 &\le\sum_{j,k}\Delta s_j
 \int_{t_{k-1}}^{t_k}p(f(\xi_j,t)-f(\xi_j,\eta_k))\,\mathrm dt\\
 &\le(b-a)(d-c)\omega_p^{(2)}(|Q|).
\end{split}
\]
Istotnie, na każdym podprzedziale całkujemy różnicę
od funkcji stałej o wartości $f(\xi_j,\eta_k)$, po czym
stosujemy addytywność i oszacowanie seminormowe.
Dla $B(t)=\int_a^b f(s,t)\,\mathrm ds$ analogicznie
$p(S(B,Q)-D(P,Q))\le(b-a)(d-c)\omega_p^{(1)}(|P|)$.
Wybierzmy ciągi podziałów $P_n,Q_n$ o średnicach dążących
do zera. Sumy $S(A,P_n)$ i $S(B,Q_n)$ zbiegają do
odpowiednich całek iterowanych z twierdzenia o istnieniu
i ciągłości $A,B$. Powyższe oszacowania dają
$p(S(A,P_n)-S(B,Q_n))\to0$ dla każdej $p$.
Ciągłość seminorm i rozdzielanie punktów dają równość
całek iterowanych.
\end{proof}

\subsection{Przykłady funkcyjne i całkowanie rodzin operatorów}

\begin{przyklad}[Całkowanie w $C^\infty(\Omega)$]
\label{ex:c-infty-integral15a}
Niech $d\ge1$, $\K\in\{\R,\C\}$, $a<b$ będą rzeczywiste
i $\Omega\subset\R^d$ będzie niepustym zbiorem otwartym.
Przestrzeń $C^\infty(\Omega;\K)$ funkcji skalarnych
wyposażamy w seminormy
$p_{K,m}(u)=\max_{|\alpha|\le m}\sup_{x\in K}|\partial^\alpha u(x)|$,
gdzie $K\subseteq\Omega$ jest niepustym zbiorem zwartym,
a $m\in\N_0$.
Jest to zupełna przestrzeń lokalnie wypukła, co szczegółowo
udowodnimy w \hyperref[sec:gladkie-schwartz17]{rozdziale~17}. Jeżeli $t\mapsto u_t$ jest
ciągłe z $[a,b]$ do tej przestrzeni, to
$u=\int_a^b u_t\,\mathrm dt\in C^\infty(\Omega)$ i
\[
 \partial^\alpha u(x)
 =\int_a^b\partial^\alpha u_t(x)\,\mathrm dt
 \qquad(x\in\Omega,\ \alpha\in\N_0^d).
\]
Istotnie, $p_{K,m}(\partial^\alpha v)\le p_{K,m+|\alpha|}(v)$,
więc $\partial^\alpha$ jest ciągłym operatorem liniowym.
Ewaluacja $v\mapsto v(x)$ jest ciągła, ponieważ
$|v(x)|=p_{\{x\},0}(v)$. Z dwukrotnego zastosowania
przemienności całki z odwzorowaniem liniowym ciągłym
otrzymujemy wskazany wzór, bez zakładania różniczkowania
pod całką skalarną.
\end{przyklad}

\begin{przyklad}[Całkowanie w przestrzeni Schwartza]
\label{ex:schwartz-integral15a}
Niech $d\ge1$, $\K\in\{\R,\C\}$ i $a<b$ będą rzeczywiste.
Niech $\mathcal S(\R^d)=\mathcal S(\R^d;\K)$ oznacza
przestrzeń funkcji gładkich o wartościach w $\K$,
dla których $q_{m,k}(u)=\max_{|\alpha|\le k}\sup_{x\in\R^d}
(1+|x|)^m|\partial^\alpha u(x)|$ jest skończone
przy wszystkich $m,k\in\N_0$; tę topologię również
omówimy w \hyperref[sec:gladkie-schwartz17]{rozdziale~17}. Przestrzeń ta jest zupełna.
Dla ciągłego $t\mapsto u_t\in\mathcal S(\R^d)$ całka
$u=\int_a^b u_t\,\mathrm dt$ należy do $\mathcal S(\R^d)$
i spełnia $q_{m,k}(u)\le\int_a^b q_{m,k}(u_t)\,\mathrm dt$.
Różniczkowanie jest na tej przestrzeni ciągłym
odwzorowaniem liniowym: dla każdego wielowskaźnika $\alpha$
seminorma $q_{m,k}(\partial^\alpha u)$ jest ograniczona
przez $q_{m,k+|\alpha|}(u)$. Stąd
$\partial^\alpha(\int_a^b u_t\,\mathrm dt)
=\int_a^b\partial^\alpha u_t\,\mathrm dt$
(jako równość w przestrzeni Schwartza).
\end{przyklad}

\begin{twierdzenie}[Silna całka rodziny operatorów]
\label{thm:strong-operator-integral15a}
Niech $a<b$ będą rzeczywiste, $X,Y$ przestrzeniami Banacha
nad wspólnym ciałem $\K\in\{\R,\C\}$ i
$T:[a,b]\to\mathcal L(X,Y)$ rodziną operatorów taką, że
$t\mapsto T(t)x$ jest ciągła w $Y$ dla każdego $x\in X$.
Wtedy wzór
$Ax=\int_a^b T(t)x\,\mathrm dt$ określa operator
$A\in\mathcal L(X,Y)$, a
\[
 \|A\|\le(b-a)\sup_{t\in[a,b]}\|T(t)\|<\infty.
\]
Nie wynika stąd ciągłość $t\mapsto T(t)$ w normie operatorowej.
\end{twierdzenie}
\begin{proof}
Dla każdego $x$ zbiór $\{T(t)x:t\in[a,b]\}$ jest ograniczony,
gdyż jest ciągłym obrazem zwartego odcinka. Twierdzenie
Banacha--Steinhausa zapewnia
$M=\sup_{t\in[a,b]}\|T(t)\|<\infty$.
Istnienie całki wektorowej w przestrzeni Banacha $Y$
wynika z twierdzenia o istnieniu, a liniowość z
liniowości całki. Dla $x\in X$ oszacowanie seminormowe
daje $\|Ax\|\le\int_a^b\|T(t)x\|\,\mathrm dt
\le(b-a)M\|x\|$, co dowodzi ciągłości $A$.
\end{proof}

\begin{przyklad}[Ciągłość silna nie oznacza ciągłości normowej]
\label{ex:strong-not-norm-cont15a}
W zespolonej przestrzeni Hilberta $H=L^2(\R,\lambda_1;\C)$ niech
$(U_t f)(x)=e^{itx}f(x)$ dla $f\in H$ i $t\in\R$.
Każdy $U_t$ jest liniową izometrią, ponieważ $|e^{itx}|=1$,
i $U_{t+s}=U_tU_s$, $U_0=I_H$. Dla każdego $f\in H$
i $t\to0$ mamy
$\|U_tf-f\|_2^2=\int_\R |e^{itx}-1|^2|f(x)|^2\,\mathrm d\lambda_1(x)
\to0$ z twierdzenia o zbieżności zmajoryzowanej, gdyż
funkcja podcałkowa jest ograniczona przez $4|f|^2$.
Dla dowolnego $s\in\R$ równość
$\|U_{s+t}f-U_sf\|_2=\|U_tf-f\|_2$ daje silną ciągłość w $s$.
Jednak dla $t\ne0$ mamy $\|U_t-I_H\|\le2$ z nierówności
trójkąta. Dla dowolnego $0<\varepsilon<2$ wybieramy
ograniczony otwarty przedział $A$ zawierający $\pi/t$,
na którym $|e^{itx}-1|>2-\varepsilon$; istnieje on z
ciągłości mnożnika i jego wartości w $\pi/t$.
Wektor $v=\chi_A/\sqrt{\lambda_1(A)}$ ma normę $1$, a
$\|(U_t-I_H)v\|_2^2\ge(2-\varepsilon)^2$.
Przechodząc z $\varepsilon\downarrow0$, otrzymujemy
$\|U_t-I_H\|=2$ dla każdego $t\ne0$.
Można więc całkować wektory $U_tf$ względem $t$ w $H$,
ale nie wolno bez dodatkowego dowodu zakładać ciągłości
$t\mapsto U_t$ w normie $\mathcal L(H)$.
\end{przyklad}

\begin{uwaga}[Granice zakresu]
Definicja sum Riemanna wymaga jedynie topologii i działań liniowych,
ale powyższe twierdzenia o istnieniu i oszacowaniach wykorzystują
lokalną wypukłość, seminormy i zupełność ciągową. W przestrzeni
niezupełnej ciągłość nie gwarantuje istnienia całki; w ogólnej
przestrzeni liniowo-topologicznej nielokalnie wypukłej nie można
bez uzasadnienia stosować oszacowań przez rodzinę seminorm.
Dla funkcji mierzalnych konieczne są dodatkowe pojęcia mierzalności;
teorię tę rozwiniemy osobno w \hyperref[sec:integral-vector15b]{module~15B}.
\end{uwaga}

\addtocounter{section}{-1}
\endgroup
